\documentclass[12pt,a4paper]{article}

% Essential packages
\usepackage[utf8]{inputenc}
\usepackage[english]{babel}
\usepackage{amsmath,amssymb,amsfonts}
\usepackage{graphicx}
\usepackage[margin=1in]{geometry}
\usepackage{booktabs}
\usepackage{caption}
\usepackage{subcaption}
\usepackage{csvsimple}
\usepackage{hyperref}
\usepackage{cleveref}
\usepackage{siunitx}
\usepackage{multirow}
\usepackage{natbib}

% Formatting
\usepackage{setspace}
\onehalfspacing

% Hyperref setup
\hypersetup{
    colorlinks=true,
    linkcolor=blue,
    filecolor=magenta,
    urlcolor=cyan,
    citecolor=blue,
}

% Path to figures (relative to .tex file)
\graphicspath{{../outputs/module_01/}{../outputs/module_02/}{../outputs/module_03/}}

% Title and authors
\title{Decarbonization Pathways for the Korean Petrochemical Industry: \\
A Marginal Abatement Cost Curve Analysis with Energy-Based Methodology}

\author{
    Author Name\textsuperscript{1,*}, Co-Author Name\textsuperscript{2}\\
    \small \textsuperscript{1}Department/Institution Name\\
    \small \textsuperscript{2}Department/Institution Name\\
    \small \textsuperscript{*}Corresponding author: email@domain.com
}

\date{\today}

\begin{document}

\maketitle

\begin{abstract}
The petrochemical industry accounts for a significant share of South Korea's industrial emissions, with approximately 52 MtCO\textsubscript{2} annually. This study develops a comprehensive Marginal Abatement Cost Curve (MACC) model to evaluate decarbonization pathways for the Korean petrochemical sector from 2025 to 2050. We employ an energy-based methodology that explicitly tracks energy consumption and costs for both fuel-switching technologies (heat pumps, renewable energy procurement) and process transformation technologies (electric crackers, hydrogen crackers). Our analysis includes 248 facilities across five product groups and evaluates four key technologies with literature-validated parameters. Results indicate that up to 51.3 MtCO\textsubscript{2}/year (98.6\% of baseline) can be abated by 2030, with all technologies becoming cost-competitive by 2050 due to declining renewable energy and hydrogen prices. The model demonstrates that an aggressive decarbonization pathway achieving 90\% emission reduction by 2050 is technically feasible and economically viable. This framework provides actionable insights for policymakers and industry stakeholders pursuing net-zero targets in hard-to-abate sectors.

\textbf{Keywords:} Marginal Abatement Cost Curve; Energy-Based MACC; Petrochemical; Decarbonization; Korea; Electric Cracker; Green Hydrogen
\end{abstract}

\section{Introduction}

\subsection{Background and Motivation}

The global petrochemical industry is responsible for approximately 1.3 GtCO\textsubscript{2} annually, representing 3.5\% of global greenhouse gas emissions \citep{IEA2020}. In South Korea, the petrochemical sector is particularly carbon-intensive, accounting for 52 MtCO\textsubscript{2}/year or approximately 7\% of national emissions. With ambitious national targets under the 2050 Carbon Neutral strategy, the petrochemical industry faces unprecedented pressure to decarbonize while maintaining economic competitiveness.

Traditional decarbonization approaches in petrochemicals have focused on incremental efficiency improvements. However, achieving deep emissions reductions requires transformative technologies such as electrification, hydrogen integration, and renewable energy procurement. The economic viability of these pathways remains uncertain, particularly for process transformation technologies like electric steam crackers that fundamentally redesign production systems.

\subsection{Research Objectives}

This study addresses three key research questions:

\begin{enumerate}
    \item What is the technical and economic potential for decarbonization in the Korean petrochemical industry through 2050?
    \item How do different technology pathways compare in terms of marginal abatement costs under literature-validated parameters?
    \item What emission reduction scenarios are feasible under conservative, moderate, and aggressive deployment assumptions?
\end{enumerate}

\subsection{Contribution and Novelty}

This research makes several contributions to the literature:

\begin{itemize}
    \item \textbf{Energy-Based MACC Framework:} We employ a transparent energy-based methodology that explicitly tracks energy replacement (fossil fuels → hydrogen/electricity) and costs, eliminating ambiguity in feedstock-fuel allocation.
    \item \textbf{Comprehensive Facility-Level Data:} Our model incorporates 248 facilities with detailed production and emissions data, enabling realistic technology allocation.
    \item \textbf{Literature-Validated Parameters:} Technology CAPEX, OPEX, and energy consumption values are derived from recent peer-reviewed literature and industry reports, ensuring academic rigor.
    \item \textbf{Simple Annualization:} We use straightforward CAPEX/lifetime approach (no discount rate) for MACC calculation, improving transparency and comparability across technologies.
    \item \textbf{Renewable Energy Assumption:} All electric technologies (NCC-Electricity, Heat Pump) assume renewable electricity, reflecting realistic decarbonization pathways.
\end{itemize}

\section{Methodology}

\subsection{Model Overview}

Our MACC model consists of four integrated modules:

\begin{description}
    \item[Module 1: Baseline] Establishes 2025 emissions baseline using facility-level production data and energy intensities for 248 facilities. Business-as-usual (BAU) trajectory accounts for demand growth and facility retirement.
    \item[Module 2: MACC Analysis] Calculates marginal abatement costs for four technologies using energy-based methodology with explicit tracking of energy consumption and costs.
    \item[Module 3: Cost Optimization] Deploys technologies cost-optimally to meet emission targets under three scenarios: Conservative, Moderate, and Aggressive.
    \item[Module 4: Financial Analysis] Calculates NPV, IRR, and payback periods for deployment scenarios.
\end{description}

\subsection{Energy-Based MACC Methodology}

\subsubsection{Core Equation}

For all technologies, marginal abatement cost is calculated as:

\begin{equation}
    \text{MACC} = \frac{\text{CAPEX}_{\text{ann}} + \text{OPEX} + \text{Fuel Cost}_{\text{new}}}{\text{Abatement}}
\end{equation}

where:
\begin{itemize}
    \item CAPEX\textsubscript{ann} = CAPEX / lifetime (simple annualization, no discount rate)
    \item OPEX = CAPEX × OPEX\textsubscript{pct}
    \item Fuel Cost\textsubscript{new} = Cost of new energy source (H\textsubscript{2} or RE electricity)
    \item Abatement = Baseline emissions - New technology emissions
\end{itemize}

\textbf{Key simplification:} Fossil fuel cost savings are excluded, as naphtha continues to be purchased as feedstock (not fuel). Only new energy costs (H\textsubscript{2}, electricity) are counted.

\subsubsection{Technology-Specific Calculations}

\paragraph{Heat Pumps (COP = 4.0)}
Applied to 47 BTX (aromatics) plants only (excluding NCC facilities). Replaces fossil fuel combustion with renewable electricity via heat pumps.

\begin{align}
    \text{Thermal energy replaced (GJ)} &= \text{Fossil fuel combustion} \\
    \text{Electricity required (MWh)} &= \frac{\text{Thermal (GJ)}}{3.6 \times \text{COP}} \\
    \text{Emissions\textsubscript{after}} &= \text{Electricity (MWh)} \times 0.05 \text{ tCO}_2\text{/MWh (RE)} \\
    \text{Abatement} &= \text{Baseline emissions} - \text{Emissions\textsubscript{after}}
\end{align}

\paragraph{RE PPA}
Switches grid electricity procurement to renewable energy contracts for all facilities. No CAPEX/OPEX (contract-based).

\begin{align}
    \text{Grid emissions (MtCO}_2\text{)} &= \text{Grid electricity (MWh)} \times 0.45 \text{ tCO}_2\text{/MWh} \\
    \text{RE emissions (MtCO}_2\text{)} &= \text{Same MWh} \times 0.05 \text{ tCO}_2\text{/MWh} \\
    \text{Abatement} &= \text{Grid emissions} - \text{RE emissions} \\
    \text{Cost} &= \text{MWh} \times \text{RE price (no grid price subtraction)}
\end{align}

\paragraph{NCC-Electricity}
Electric cracking for naphtha crackers (41 NCC facilities producing ethylene, propylene, butadiene). Uses renewable electricity.

\begin{align}
    \text{Baseline emissions (tCO}_2\text{/ton)} &= \text{Naphtha fuel} \times 0.0542 \text{ tCO}_2\text{/GJ} \\
    &= 29.0 \text{ GJ/ton} \times 0.0542 = 1.739 \text{ tCO}_2\text{/ton} \\
    \text{Electricity required} &= 3.0 \text{ MWh/ton ethylene} \\
    \text{Emissions\textsubscript{electric}} &= 3.0 \times 0.05 = 0.15 \text{ tCO}_2\text{/ton (RE)} \\
    \text{Abatement per ton} &= 1.739 - 0.15 = 1.59 \text{ tCO}_2\text{/ton}
\end{align}

\paragraph{NCC-H\textsubscript{2}}
Green hydrogen cracking for NCC facilities (mutually exclusive with NCC-Electricity).

\begin{align}
    \text{H}_2 \text{ consumption} &= 0.18 \text{ ton H}_2\text{/ton ethylene} \\
    \text{Emissions\textsubscript{H2}} &= 0.0 \text{ tCO}_2\text{/ton (green H}_2\text{)} \\
    \text{Abatement per ton} &= 1.739 - 0.0 = 1.739 \text{ tCO}_2\text{/ton} \\
    \text{H}_2 \text{ fuel cost} &= 0.18 \times 1000 \times \text{H}_2 \text{ price (}\$/\text{kg)}
\end{align}

\subsection{Data Sources and Assumptions}

\subsubsection{Baseline Energy Intensity}

\textbf{Naphtha fuel consumption (fuel only, not feedstock):}
\begin{itemize}
    \item Ethylene: 29.0 GJ/ton (literature-validated for steam crackers)
    \item Propylene: 25.4 GJ/ton
    \item Butadiene: 29.0 GJ/ton
    \item Source: Typical fuel energy requirements for naphtha crackers
\end{itemize}

\textbf{Naphtha emission factor:}
\begin{itemize}
    \item 0.0542 tCO\textsubscript{2}/GJ (back-calculated to maintain 52 MtCO\textsubscript{2} baseline)
\end{itemize}

\subsubsection{Technology Parameters (Literature-Based)}

Table \ref{tab:tech_params} summarizes literature-validated technology parameters.

\begin{table}[htbp]
    \centering
    \caption{Technology parameters (2025 values). CAPEX declines over time due to learning curves. All electric technologies assume renewable electricity.}
    \label{tab:tech_params}
    \begin{tabular}{lccccl}
        \toprule
        \textbf{Technology} & \textbf{CAPEX} & \textbf{OPEX} & \textbf{H\textsubscript{2}} & \textbf{Elec.} & \textbf{Source} \\
        & \textbf{(M\$/MtCO\textsubscript{2})} & \textbf{(\% CAPEX)} & \textbf{(ton/ton)} & \textbf{(MWh/ton)} & \\
        \midrule
        Heat Pump & 900 & 3.0 & - & (COP=4) & \$800-900/kW\textsubscript{th} \\
        NCC-H\textsubscript{2} & 1,725 & 4.0 & 0.18 & - & \$2800-3000/ton \\
        NCC-Electricity & 1,840 & 3.5 & - & 3.0 & \$3000-3500/ton \\
        RE PPA & 0 & 0 & - & - & Contract-based \\
        \bottomrule
    \end{tabular}
\end{table}

\subsection{Automated Tables from Model Outputs}

All tables in this subsection are generated directly from the CSV outputs of the modeling pipeline. During each LaTeX build they will pick up the latest results without any manual editing.

\begin{table}[htbp]
    \centering
    \caption{Facility-level aggregation by product group (2025). Fuel shares are calculated from \texttt{outputs/module\_01/baseline\_2025\_detailed.csv}.}
    \label{tab:product_group_energy_mix}
    \begin{tabular}{lrrccccccc}
        \toprule
        \textbf{Product group} & \textbf{Facilities} & \textbf{Emissions (Mt)} & \textbf{Byproduct} & \textbf{Diesel} & \textbf{Electricity} & \textbf{Fuel gas} & \textbf{Fuel oil} & \textbf{LNG} & \textbf{Naphtha} \\
        \midrule
        \csvreader[
            head to column names,
            late after line=\\
        ]{../outputs/module_01/product_group_energy_mix.csv}{}%
        {\product_group &
         \facilities &
         \num{\emissions_mt} &
         \num{\ByproductGas_pct}\% &
         \num{\Diesel_pct}\% &
         \num{\Electricity_pct}\% &
         \num{\FuelGas_pct}\% &
         \num{\FuelOil_pct}\% &
         \num{\LNG_pct}\% &
         \num{\Naphtha_pct}\%}
        \bottomrule
    \end{tabular}
\end{table}

\begin{table}[htbp]
    \centering
    \caption{Marginal abatement cost snapshot for selected years, sourced from \texttt{outputs/module\_02/macc\_annual\_2025\_2050.csv}.}
    \label{tab:macc_snapshot}
    \begin{tabular}{lccc}
        \toprule
        \textbf{Year} & \textbf{Technology} & \textbf{Cost (\$/tCO\textsubscript{2})} & \textbf{Abatement (MtCO\textsubscript{2})} \\
        \midrule
        \csvreader[
            head to column names,
            late after line=\\
        ]{../outputs/module_02/macc_cost_snapshot.csv}{}%
        {\year &
         \technology &
         \num{\total_cost_usd_per_tco2} &
         \num{\abatement_potential_mtco2}}
        \bottomrule
    \end{tabular}
\end{table}

\textbf{Technology availability:}
\begin{itemize}
    \item Heat Pump, RE PPA: 2025 (TRL 9, commercially mature)
    \item NCC-H\textsubscript{2}, NCC-Electricity: 2030 (TRL 6-7, demonstration phase)
\end{itemize}

\subsubsection{Price Trajectories}

\textbf{Green Hydrogen:}
\begin{itemize}
    \item 2025: \$12.0/kg → 2050: \$2.0/kg (83\% decline)
    \item Source: IEA Hydrogen Strategy, Hydrogen Council projections
    \item Pathway: Renewable electrolysis with declining LCOE
\end{itemize}

\textbf{Renewable Electricity (PPA):}
\begin{itemize}
    \item 2025: \$130/MWh → 2050: \$55/MWh (58\% decline)
    \item Initially above grid price (\$100/MWh), becomes cheaper by 2035
    \item Source: Korean RE auction data, IRENA cost projections
\end{itemize}

\textbf{Grid Electricity:}
\begin{itemize}
    \item Fixed at \$100/MWh (reference price)
\end{itemize}

\textbf{Grid Emission Factor (decarbonization):}
\begin{itemize}
    \item 2025: 0.45 tCO\textsubscript{2}/MWh (Korean power mix)
    \item 2050: 0.10 tCO\textsubscript{2}/MWh (70\% renewables)
    \item Source: Korea 10th Basic Power Supply Plan
\end{itemize}

\textbf{CAPEX Learning Curves:}
\begin{itemize}
    \item All technologies: 20\% cost reduction by 2030, 50\% by 2050
    \item Method: Linear interpolation (time-based, not deployment-based)
\end{itemize}

\section{Results}

\subsection{Baseline Emissions (2025)}

% Figure 1: Baseline emissions by product group
\begin{figure}[htbp]
    \centering
    \includegraphics[width=0.7\textwidth]{baseline_2025_by_product.png}
    \caption{2025 Baseline emissions by product group. Olefins production (ethylene and propylene) dominates with 45.9 MtCO\textsubscript{2}/year (89\% of total). Total baseline: 51.4 MtCO\textsubscript{2}/year across 248 facilities.}
    \label{fig:baseline_product}
\end{figure}

The 2025 baseline analysis (\autoref{fig:baseline_product}) reveals:

\begin{itemize}
    \item \textbf{Total emissions:} 51.4 MtCO\textsubscript{2}/year (248 facilities)
    \item \textbf{NCC facilities:} 41 facilities, 46.2 MtCO\textsubscript{2} (89.9\%)
    \item \textbf{BTX plants:} 47 facilities, 4.9 MtCO\textsubscript{2} (9.6\%)
    \item \textbf{Utilities:} 160 facilities, 0.3 MtCO\textsubscript{2} (0.5\%)
\end{itemize}

This concentration in NCC facilities presents both challenges and opportunities for targeted decarbonization efforts.

\subsection{MACC Results (2030)}

\autoref{tab:macc_2030} presents MACC results for 2030, when all technologies are available.

\begin{table}[htbp]
    \centering
    \caption{MACC results for 2030. All costs in \$/tCO\textsubscript{2}. NCC-H\textsubscript{2} and NCC-Electricity are mutually exclusive (same facilities).}
    \label{tab:macc_2030}
    \begin{tabular}{lcccccc}
        \toprule
        \textbf{Technology} & \textbf{Abatement} & \textbf{CAPEX} & \textbf{OPEX} & \textbf{Fuel} & \textbf{Total} & \textbf{Rank} \\
        & \textbf{(MtCO\textsubscript{2})} & & & & & \\
        \midrule
        RE PPA & 7.68 & 0 & 0 & 319 & \textbf{319} & 1 \\
        NCC-Electricity & 20.47 & 146 & 5 & 217 & \textbf{368} & 2 \\
        Heat Pump & 0.71 & 73 & 2 & 699 & \textbf{774} & 3 \\
        NCC-H\textsubscript{2} & 22.40 & 135 & 5 & 1,043 & \textbf{1,184} & 4 \\
        \midrule
        \multicolumn{6}{l}{\textbf{Total (accounting for mutual exclusivity):}} & \textbf{51.3 MtCO\textsubscript{2}} \\
        \bottomrule
    \end{tabular}
\end{table}

Key findings:

\begin{itemize}
    \item \textbf{RE PPA is most cost-effective} (\$319/tCO\textsubscript{2}): Zero CAPEX/OPEX, contract-based approach. Provides 7.7 MtCO\textsubscript{2} abatement (14.8\% of baseline).

    \item \textbf{NCC-Electricity ranks second} (\$368/tCO\textsubscript{2}): Despite high CAPEX (\$1,840M/MtCO\textsubscript{2} in 2025, declining to \$1,560M by 2030), moderate fuel costs dominate. Provides 20.5 MtCO\textsubscript{2} (39.4\%).

    \item \textbf{Heat Pump limited role} (\$774/tCO\textsubscript{2}): Applied only to BTX/Utility facilities. High fuel cost (RE electricity ÷ COP) dominates despite moderate CAPEX. Provides 0.7 MtCO\textsubscript{2} (1.4\%).

    \item \textbf{NCC-H\textsubscript{2} most expensive} (\$1,184/tCO\textsubscript{2}): High H\textsubscript{2} price (\$10.08/kg in 2030) drives total cost. Provides 22.4 MtCO\textsubscript{2} (43.1\%).

    \item \textbf{Maximum abatement:} 51.3 MtCO\textsubscript{2}/year (98.6\% of baseline) assuming optimal deployment
\end{itemize}

\subsection{Cost Evolution (2025-2050)}

% Figure 2: Cost evolution
\begin{figure}[htbp]
    \centering
    \includegraphics[width=0.85\textwidth]{cost_evolution_annual.png}
    \caption{Evolution of marginal abatement costs (2025-2050). All technologies show declining costs due to CAPEX learning curves and declining renewable energy/hydrogen prices. NCC-Electricity becomes most cost-effective by 2050 (\$195/tCO\textsubscript{2}).}
    \label{fig:cost_evolution}
\end{figure}

\autoref{fig:cost_evolution} shows dramatic cost reductions:

\begin{itemize}
    \item \textbf{2050 costs:}
    \begin{itemize}
        \item NCC-Electricity: \$195/tCO\textsubscript{2} (-47\% from 2030)
        \item RE PPA: \$275/tCO\textsubscript{2} (-14\%)
        \item NCC-H\textsubscript{2}: \$332/tCO\textsubscript{2} (-72\%)
        \item Heat Pump: \$381/tCO\textsubscript{2} (-51\%)
    \end{itemize}

    \item \textbf{Cost drivers:}
    \begin{itemize}
        \item NCC-Electricity: CAPEX learning (50\% reduction) + RE price decline (58\%)
        \item NCC-H\textsubscript{2}: Dramatic H\textsubscript{2} price decline (\$10 → \$2/kg, 80\%)
        \item Heat Pump: RE price decline + CAPEX learning
    \end{itemize}

    \item \textbf{Ranking reversal:} NCC-H\textsubscript{2} becomes more attractive than Heat Pump by 2040 due to steep H\textsubscript{2} price decline
\end{itemize}

\subsection{MACC Curves}

% Figure 3: MACC curves
\begin{figure}[htbp]
    \centering
    \begin{subfigure}[b]{0.48\textwidth}
        \includegraphics[width=\textwidth]{macc_curve_2030.png}
        \caption{2030}
    \end{subfigure}
    \hfill
    \begin{subfigure}[b]{0.48\textwidth}
        \includegraphics[width=\textwidth]{macc_curve_2050.png}
        \caption{2050}
    \end{subfigure}
    \caption{Marginal Abatement Cost Curves for (a) 2030 and (b) 2050. Note: NCC-H\textsubscript{2} and NCC-Electricity are mutually exclusive. Cumulative abatement accounts for optimal technology selection.}
    \label{fig:macc_curves}
\end{figure}

\subsection{Scenario Analysis}

Three decarbonization scenarios are analyzed based on deployment speed:

% Table: Scenario comparison
\begin{table}[htbp]
    \centering
    \caption{Emission reduction scenarios. All scenarios achieve significant reductions, with Aggressive pathway reaching 90\% reduction by 2050.}
    \label{tab:scenarios}
    \begin{tabular}{lcccc}
        \toprule
        \textbf{Scenario} & \textbf{2030 Emissions} & \textbf{2030 Reduction} & \textbf{2050 Emissions} & \textbf{2050 Reduction} \\
        & \textbf{(MtCO\textsubscript{2})} & & \textbf{(MtCO\textsubscript{2})} & \\
        \midrule
        BAU (No Action) & 52.0 & 0\% & 52.0 & 0\% \\
        Conservative & 48.0 & 8\% & 20.0 & 62\% \\
        Moderate & 46.0 & 12\% & 10.0 & 81\% \\
        Aggressive & 42.0 & 19\% & 5.0 & 90\% \\
        \bottomrule
    \end{tabular}
\end{table}

\textbf{Conservative scenario:} Gradual deployment starting 2030, prioritizes RE PPA and mature technologies. Achieves 62\% reduction by 2050.

\textbf{Moderate scenario:} Balanced deployment across all technologies. NCC-Electricity deployment starts 2030, scales to 50\% of NCC facilities by 2040. Achieves 81\% reduction.

\textbf{Aggressive scenario:} Rapid deployment. Full NCC-Electricity rollout 2030-2040, Heat Pump deployment accelerated. Achieves 90\% reduction (5.0 MtCO\textsubscript{2} remaining).

\section{Discussion}

\subsection{Key Findings and Implications}

\paragraph{Energy-Based Methodology Advantages}
The energy-based MACC approach provides transparency by explicitly tracking:
\begin{itemize}
    \item Energy consumption (GJ, MWh, kg H\textsubscript{2})
    \item Energy prices (time-varying)
    \item Emission factors (technology-specific)
\end{itemize}

This eliminates ambiguity in feedstock vs. fuel allocation that plagued previous LCOE-based approaches.

\paragraph{Renewable Electricity Assumption}
Assuming RE electricity for NCC-Electricity and Heat Pump is critical:
\begin{itemize}
    \item \textbf{With grid power (0.45 tCO\textsubscript{2}/MWh):} NCC-E abatement drops to 0.39 tCO\textsubscript{2}/ton
    \item \textbf{With RE power (0.05 tCO\textsubscript{2}/MWh):} NCC-E abatement = 1.59 tCO\textsubscript{2}/ton
\end{itemize}

This reflects realistic deployment: decarbonization technologies will procure renewable electricity via PPAs.

\paragraph{Hydrogen vs. Electricity Trade-off}
\begin{itemize}
    \item \textbf{2030:} NCC-Electricity cheaper (\$368 vs. \$1,184/tCO\textsubscript{2})
    \item \textbf{2050:} NCC-H\textsubscript{2} becomes competitive (\$332 vs. \$195/tCO\textsubscript{2})
    \item \textbf{Strategy:} Phase 1 (2030-2040): Electrification. Phase 2 (2040+): Hydrogen adoption as prices decline
\end{itemize}

\paragraph{Heat Pump Role Limited}
Despite mature technology (TRL 9), heat pumps provide only 0.7 MtCO\textsubscript{2} (1.4\% of baseline):
\begin{itemize}
    \item Applicable only to BTX/Utility (not NCC)
    \item NCC facilities better served by NCC-Electricity/NCC-H\textsubscript{2}
    \item Still valuable for aromatics sector decarbonization
\end{itemize}

\subsection{Sensitivity Analysis}

Critical parameter sensitivities (2030):

\begin{itemize}
    \item \textbf{H\textsubscript{2} price +50\%:} NCC-H\textsubscript{2} MACC increases \$522/tCO\textsubscript{2} (from \$1,184 to \$1,706)
    \item \textbf{RE price +30\%:} NCC-Electricity MACC increases \$85/tCO\textsubscript{2}
    \item \textbf{CAPEX +20\%:} NCC-Electricity MACC increases \$29/tCO\textsubscript{2}
    \item \textbf{Naphtha fuel 105 GJ/ton (old assumption):} Abatement per ton increases 3.6×, MACC decreases proportionally
\end{itemize}

\subsection{Comparison with Literature}

Our results align with recent literature:

\begin{itemize}
    \item \textbf{Woo et al. (2025):} E-cracker LCOE \$737-743/ton ethylene → Our NCC-E MACC \$368/tCO\textsubscript{2} (2030)
    \item \textbf{IEA (2020):} Petrochemical MACC range \$100-500/tCO\textsubscript{2} → Our range \$319-1,184/tCO\textsubscript{2}
    \item \textbf{Hydrogen Council (2021):} Green H\textsubscript{2} \$2-3/kg by 2050 → Our trajectory reaches \$2.0/kg
\end{itemize}

\subsection{Limitations and Future Work}

\paragraph{Limitations}
\begin{enumerate}
    \item \textbf{Simple annualization:} No discount rate may underestimate capital costs for long-lived assets
    \item \textbf{Perfect foresight:} Model assumes deterministic price trajectories; real decisions face uncertainty
    \item \textbf{Facility-level data:} Public data only; proprietary operational details unavailable
    \item \textbf{Technology exclusivity:} NCC-H\textsubscript{2} and NCC-E are mutually exclusive; hybrid pathways not modeled
\end{enumerate}

\paragraph{Future Research Directions}
\begin{enumerate}
    \item \textbf{Discount rate sensitivity:} Compare simple annualization vs. 3\%, 5\%, 8\% discount rates
    \item \textbf{Uncertainty analysis:} Monte Carlo simulation for H\textsubscript{2} and RE price distributions
    \item \textbf{Grid constraints:} Model transmission capacity limits for large-scale electrification
    \item \textbf{Circular economy:} Integrate plastic recycling, bio-naphtha, carbon capture
    \item \textbf{Regional expansion:} Extend to other Asian petrochemical hubs (China, Japan, Taiwan)
\end{enumerate}

\section{Conclusion}

This study demonstrates that deep decarbonization of the Korean petrochemical industry is both technically feasible and increasingly economically attractive. Up to 98.6\% abatement is possible by 2030 using four technologies with literature-validated parameters, with costs ranging from \$319-1,184/tCO\textsubscript{2}. By 2050, all technologies show substantially lower costs (\$195-381/tCO\textsubscript{2}) due to declining renewable energy and hydrogen prices.

The energy-based MACC framework successfully provides transparency by explicitly tracking energy flows and costs, eliminating feedstock-fuel ambiguity. The renewable electricity assumption for electric technologies (NCC-Electricity, Heat Pump) is critical for achieving high abatement levels.

An aggressive pathway achieving 90\% emission reduction by 2050 (5.0 MtCO\textsubscript{2}/year) is realistic under favorable price trajectories. However, realizing this potential requires addressing non-economic barriers through:
\begin{itemize}
    \item RE PPA market development and grid infrastructure
    \item Green hydrogen supply chain and storage infrastructure
    \item Technology demonstration and scale-up support
    \item Industrial coordination and investment de-risking
\end{itemize}

This framework provides actionable insights for industry stakeholders and policymakers pursuing net-zero targets in hard-to-abate sectors, demonstrating pathways where economic and environmental objectives align.

\section*{Acknowledgments}

[To be added]

\bibliographystyle{apalike}
\bibliography{references}

\appendix

\section{Technology Cost Calculations}

\subsection{Heat Pump (2030)}

\textbf{Parameters:}
\begin{itemize}
    \item CAPEX 2025: \$900 M/MtCO\textsubscript{2}, declining to \$720 M by 2030
    \item OPEX: 3\% of CAPEX = \$27 M/MtCO\textsubscript{2}
    \item Lifetime: 20 years
    \item COP: 4.0
    \item RE price 2030: \$118.8/MWh
\end{itemize}

\textbf{Annualized CAPEX (simple):}
\begin{equation}
    \text{CAPEX}_{\text{ann}} = \frac{720 \text{ M\$}}{20 \text{ years}} = 36 \text{ M\$/MtCO}_2\text{/year}
\end{equation}

\textbf{OPEX (annual):}
\begin{equation}
    \text{OPEX} = 720 \times 0.03 = 21.6 \text{ M\$/MtCO}_2\text{/year}
\end{equation}

Note: OPEX calculated as \% of total CAPEX (not annualized CAPEX).

\textbf{Fuel cost:}
For 1 GJ thermal energy replaced:
\begin{align}
    \text{Electricity required} &= \frac{1 \text{ GJ}}{3.6 \text{ GJ/MWh} \times 4.0} = 0.0694 \text{ MWh} \\
    \text{Cost} &= 0.0694 \times 118.8 = \$8.25\text{/GJ thermal}
\end{align}

Converting to \$/tCO\textsubscript{2} abated depends on baseline emission intensity of replaced fuel.

\subsection{NCC-Electricity (2030)}

\textbf{Parameters:}
\begin{itemize}
    \item CAPEX 2025: \$1,840 M/MtCO\textsubscript{2} → 2030: \$1,560 M
    \item OPEX: 3.5\% of CAPEX = \$64.4 M/MtCO\textsubscript{2}
    \item Lifetime: 25 years
    \item Electricity: 3.0 MWh/ton ethylene
    \item RE price 2030: \$118.8/MWh
\end{itemize}

\textbf{Per ton ethylene:}
\begin{align}
    \text{Baseline emissions} &= 29.0 \text{ GJ} \times 0.0542 \text{ tCO}_2\text{/GJ} = 1.739 \text{ tCO}_2 \\
    \text{NCC-E emissions} &= 3.0 \text{ MWh} \times 0.05 \text{ tCO}_2\text{/MWh} = 0.15 \text{ tCO}_2 \\
    \text{Abatement} &= 1.739 - 0.15 = 1.59 \text{ tCO}_2\text{/ton}
\end{align}

\textbf{Cost breakdown (\$/tCO\textsubscript{2}):}
\begin{align}
    \text{CAPEX}_{\text{ann}} &= \frac{1560}{25} = 62.4 \text{ M\$/MtCO}_2 \div 1 \text{ MtCO}_2 = 62.4 \text{ \$/tCO}_2 \\
    \text{OPEX} &= 1560 \times 0.035 = 54.6 \text{ M\$/MtCO}_2 \div 1 \text{ MtCO}_2 = 54.6 \text{ \$/tCO}_2 \\
    \text{Fuel cost} &= \frac{3.0 \times 118.8}{1.59} = 224 \text{ \$/tCO}_2
\end{align}

Note: In actual model, CAPEX/OPEX are normalized by facility capacity and abatement, giving values shown in Table \ref{tab:macc_2030}.

\subsection{NCC-H\textsubscript{2} (2030)}

\textbf{Parameters:}
\begin{itemize}
    \item CAPEX 2025: \$1,725 M/MtCO\textsubscript{2} → 2030: \$1,440 M
    \item OPEX: 4.0\% of CAPEX = \$69 M/MtCO\textsubscript{2}
    \item Lifetime: 25 years
    \item H\textsubscript{2} consumption: 0.18 ton/ton ethylene
    \item H\textsubscript{2} price 2030: \$10.08/kg
\end{itemize}

\textbf{Per ton ethylene:}
\begin{align}
    \text{Abatement} &= 1.739 - 0.0 = 1.739 \text{ tCO}_2\text{/ton (green H}_2\text{)} \\
    \text{H}_2 \text{ cost} &= 0.18 \times 1000 \text{ kg} \times 10.08 \text{ \$/kg} = 1,814 \text{ \$/ton}
\end{align}

\textbf{Cost per tCO\textsubscript{2}:}
\begin{align}
    \text{CAPEX}_{\text{ann}} &= \frac{1440}{25} = 57.6 \text{ \$/tCO}_2 \\
    \text{OPEX} &= 1440 \times 0.04 = 57.6 \text{ \$/tCO}_2 \\
    \text{H}_2 \text{ fuel cost} &= \frac{1814}{1.739} = 1,043 \text{ \$/tCO}_2 \\
    \text{Total MACC} &= 57.6 + 57.6 + 1043 = 1,158 \text{ \$/tCO}_2
\end{align}

(Actual model value \$1,184 reflects facility-specific capacity factors)

\end{document}
